\documentclass[10pt,a4paper]{article}
\usepackage[utf8]{inputenc}
\usepackage[indonesian]{babel}
\usepackage{amsmath,amssymb,amsfonts}
\usepackage{geometry}
\usepackage{multicol}
\usepackage{enumitem}
\usepackage{fourier}

\geometry{top=1.5cm, bottom=1.5cm, left=1.5cm, right=1.5cm}

\pagestyle{empty}

\begin{document}

\begin{center}
  \textbf{\Large Latihan Soal Bahasa Inggris}
\end{center}
\vspace{0.3cm}

\begin{multicols}{2}
  \setlist[enumerate,1]{leftmargin=*,label=\arabic*.}
  \setlist[enumerate,2]{leftmargin=*,label=\Alph*.}
  
  \begin{enumerate}
    \item What is the definition of a narrative text?
    \begin{enumerate}
      \item A text that explains how to do something step by step
      \item A text that tells a story involving characters and a sequence of events
      \item A text that describes a particular person, place, or thing
      \item A text that argues a specific point of view
      \item A text that reports actual news events
    \end{enumerate}

    \item What is the primary purpose of a narrative text?
    \begin{enumerate}
      \item To entertain or engage the reader through a story
      \item To inform the reader about factual events
      \item To persuade the reader to buy a product
      \item To describe a scientific process
      \item To criticize a literary artwork
    \end{enumerate}

    \item Which of the following is NOT a common example of a narrative text?
    \begin{enumerate}
      \item Fairy tales
      \item Legends
      \item Fables
      \item Instruction manuals
      \item Myths
    \end{enumerate}

    \item What is the correct generic structure of a narrative text?
    \begin{enumerate}
      \item Orientation - Sequence of Events - Complication - Resolution - Coda
      \item Orientation - Complication - Sequence of Events - Resolution - Coda
      \item Complication - Orientation - Resolution - Sequence of Events - Coda
      \item Orientation - Resolution - Complication - Sequence of Events - Coda
      \item Coda - Orientation - Complication - Sequence of Events - Resolution
    \end{enumerate}

    \item Which structural element introduces the characters, setting, and time of the story?
    \begin{enumerate}
      \item Coda
      \item Resolution
      \item Sequence of Events
      \item Complication
      \item Orientation
    \end{enumerate}

    \item What is introduced in the complication stage of a narrative text?
    \begin{enumerate}
      \item The time and setting of the story
      \item The moral lesson or message
      \item The main problem or conflict faced by the characters
      \item The final outcome of the main conflict
      \item The general description of the characters
    \end{enumerate}

    \item In a narrative text, the sequence of events is usually presented in...
    \begin{enumerate}
      \item Alphabetical order
      \item Reverse chronological order
      \item Random order
      \item Chronological order
      \item Geographical order
    \end{enumerate}

    \item Which of the following statements about "Resolution" is TRUE?
    \begin{enumerate}
      \item It must always end with a happy outcome.
      \item It must always end with a tragic outcome.
      \item It can have a positive, negative, or unresolved outcome.
      \item It always states the explicit moral lesson.
      \item It introduces a completely new conflict.
    \end{enumerate}

    \item The optional closing section that states the lesson, moral, or message of the story is called...
    \begin{enumerate}
      \item Sequence
      \item Orientation
      \item Complication
      \item Resolution
      \item Coda
    \end{enumerate}

    \item Which verb tense is frequently used in narrative texts because they commonly describe events that occurred in the past?
    \begin{enumerate}
      \item Simple Present Tense
      \item Present Continuous Tense
      \item Simple Past Tense
      \item Simple Future Tense
      \item Present Perfect Tense
    \end{enumerate}

    \item Words like \textit{first, then, next, after that,} and \textit{finally} are examples of...
    \begin{enumerate}
      \item Action verbs
      \item Temporal connectives
      \item Specific participants
      \item Adjectives
      \item Pronouns
    \end{enumerate}

    \item Based on the story of "The Tortoise and the Hare", what serves as the complication?
    \begin{enumerate}
      \item A hare and a tortoise lived in a forest.
      \item The tortoise reaches the finish line first.
      \item The hare challenges the tortoise to a race.
      \item The hare becomes overconfident and takes a nap.
      \item Perseverance is more valuable than overconfidence.
    \end{enumerate}

    \item Action verbs describe events and activities performed by characters. Which of the following is an example of an action verb?
    \begin{enumerate}
      \item Beautiful
      \item Quickly
      \item Walked
      \item Very
      \item They
    \end{enumerate}

    \item Why are adjectives and adverbs often used in narrative texts (Descriptive Language)?
    \begin{enumerate}
      \item To describe characters, settings, and actions
      \item To indicate the chronological order of events
      \item To introduce the central conflict
      \item To summarize the moral lesson
      \item To replace the action verbs
    \end{enumerate}

    \item Narratives commonly refer to specific participants rather than general groups. An example of a specific participant is...
    \begin{enumerate}
      \item A general group of animals
      \item The society
      \item Belle and the Beast
      \item Humans in general
      \item People
    \end{enumerate}

    \item What is the definition of reported speech?
    \begin{enumerate}
      \item It reports the exact words of the speaker using quotation marks.
      \item It reports the meaning of what another person said without reproducing the exact words.
      \item It is a speech delivered in front of a large audience.
      \item It is a text that tells a fictional story.
      \item It only reports commands and requests.
    \end{enumerate}

    \item When the reporting verb is in the past tense, the verb in the reported clause commonly undergoes a "backshift", which means...
    \begin{enumerate}
      \item The tense is moved one step forward.
      \item The tense remains exactly the same.
      \item The tense is moved one step backward.
      \item The tense changes to future continuous.
      \item The verbs are completely omitted.
    \end{enumerate}

    \item Change this direct speech to reported speech: She said, "I like English."
    \begin{enumerate}
      \item She said that she likes English.
      \item She said that she liked English.
      \item She said that I liked English.
      \item She said that I like English.
      \item She said she had liked English.
    \end{enumerate}

    \item If the direct speech uses the Present Continuous tense (e.g., \textit{am studying}), what tense should be used in the reported speech?
    \begin{enumerate}
      \item Past Simple
      \item Present Perfect
      \item Past Continuous
      \item Past Perfect
      \item Future Simple
    \end{enumerate}

    \item Change this direct speech to reported speech: He said, "I went to Jakarta."
    \begin{enumerate}
      \item He said that he goes to Jakarta.
      \item He said that he went to Jakarta.
      \item He said that he had gone to Jakarta.
      \item He said that he has gone to Jakarta.
      \item He said that he was going to Jakarta.
    \end{enumerate}

    \item In reported speech, the time expression "tomorrow" commonly changes to...
    \begin{enumerate}
      \item that day
      \item the day before
      \item then
      \item the next day
      \item that night
    \end{enumerate}

    \item Choose the grammatically correct pattern for the verb "tell" in reported speech.
    \begin{enumerate}
      \item tell + something
      \item tell + that + statement
      \item tell + person + something
      \item tell + if + statement
      \item tell + question word
    \end{enumerate}

    \item Change to reported speech: He asked, "Are you tired?"
    \begin{enumerate}
      \item He asked me if I am tired.
      \item He asked me whether I was tired.
      \item He asked me that I was tired.
      \item He asked me why I was tired.
      \item He asked me was I tired.
    \end{enumerate}

    \item For reported WH-Questions, what is the essential rule regarding word order?
    \begin{enumerate}
      \item question word + normal question word order
      \item question word + statement word order
      \item if/whether + statement
      \item asked + to + V1
      \item that + question word + statement
    \end{enumerate}

    \item Change to reported speech: She asked, "Where do you live?"
    \begin{enumerate}
      \item She asked me where did I live.
      \item She asked me where do I live.
      \item She asked me where I lived.
      \item She asked me where I live.
      \item She asked me if I lived there.
    \end{enumerate}

    \item Commands in reported speech are commonly reported using the structure...
    \begin{enumerate}
      \item tell + person + to + V1
      \item tell + person + that + V1
      \item say + to + V1
      \item ask + if + V1
      \item said + not to + V1
    \end{enumerate}

    \item Change the negative command to reported speech: Dad said, "Don't touch that."
    \begin{enumerate}
      \item Dad told me don't touch that.
      \item Dad told me to not touch that.
      \item Dad told me not to touch that.
      \item Dad told me to touch that.
      \item Dad told me didn't touch that.
    \end{enumerate}

    \item When is a backshift in reported speech generally NOT necessary?
    \begin{enumerate}
      \item When the reporting verb is in the past tense.
      \item When the reported statement is a WH-Question.
      \item When the reported statement describes a fact that remains true.
      \item When pronouns in the sentence need to be changed.
      \item When time expressions change to a different context.
    \end{enumerate}

    \item Change to reported speech: Sarah said, "The Earth is round."
    \begin{enumerate}
      \item Sarah said that the Earth was round.
      \item Sarah said that the Earth is round.
      \item Sarah said that the Earth had been round.
      \item Sarah said if the Earth is round.
      \item Sarah said that the Earth would be round.
    \end{enumerate}

    \item Change to reported speech: David said to Maria, "I will finish this today."
    \begin{enumerate}
      \item David told Maria that he will finish this today.
      \item David told Maria that he would finish that that day.
      \item David told Maria that he would finish this that day.
      \item David told Maria that I would finish that that day.
      \item David told Maria that he had finished that that day.
    \end{enumerate}

  \end{enumerate}
\end{multicols}
\end{document}
